\documentclass[11pt]{article}
\usepackage[a4paper,margin=25mm]{geometry}
\usepackage{amsmath,amssymb}
\usepackage{hyperref}
\usepackage{url}
\setlength{\emergencystretch}{3em}
\title{Nested residual classes refute Conjecture 00000000388}
\author{gaochengzhecpu (AI-assisted submission)}
\date{3 October 2026}
\begin{document}
\maketitle

\paragraph{Statement.}
The source defines, for positive integers $m$,
\[
 W_m=\{\xi: w_n(\xi)\le n/m \text{ for infinitely many }n,
                 \quad \xi\notin U_k\text{ for every }k\}.
\]
It asserts that each $W_m$ is nonempty and uncountable, that distinct
$W_m$ are disjoint, and that their union is the set of transcendental numbers.
The nonemptiness and disjointness assertions are already incompatible with this definition.

\paragraph{Proof.}
Let $1\le m\le k$. For every nonnegative integer $n$,
\[
 \frac{n}{k}\le\frac{n}{m}.
\]
Consequently, if $w_n(\xi)\le n/k$, then $w_n(\xi)\le n/m$.
The same infinitely many indices witnessing the condition for $W_k$
therefore witness the condition for $W_m$. The exclusion from all $U_j$
is identical in both definitions. It follows that
\[
 W_k\subseteq W_m \qquad (1\le m\le k).
\]
In particular $W_2\subseteq W_1$. If the asserted nonemptiness holds,
choose $\xi\in W_2$; then $\xi\in W_1\cap W_2$, contradicting the asserted
disjointness for $1\ne2$. Equivalently, disjointness forces $W_2=\varnothing$.
This refutes the conjunction in the source, regardless of its further
uncountability and coverage assertions.\hfill$\square$

\paragraph{Scope.}
The argument does not assert the existence of any particular approximation
orders or transcendental numbers with these properties. It proves that no
assignment of ordered numerical values to the symbols $w_n(\xi)$ can make
the stated threshold definition and both assertions hold simultaneously.
It also applies if the values can be $+\infty$, or if the underlying set
is restricted to transcendental numbers. No Diophantine approximation
theorem is needed.

\newpage
\paragraph{Formal verification.}
The accompanying Lean 4.19.0 project uses only \texttt{Std}.
It proves the following statement for an arbitrary underlying type $X$,
arbitrary approximation-order data, and an arbitrary family of excluded sets $U_k$:
\[
 \neg\bigl[(\forall m\ge1,\ W_m\ne\varnothing)
      \ \wedge\ (W_m\text{ are pairwise disjoint})\bigr].
\]

\paragraph{Numerical representation and semantic bridge.}
\texttt{RationalBound} stores a numerator $a\in\mathbb N$ and a positive
denominator $b$, representing the nonnegative rational number $a/b$.
\texttt{BoundLE} compares two such fractions by cross multiplication:
$a/b\le c/d$ means $ad\le cb$. Thus this representation uses exact
rational bounds, not integer division. The theorem
\path{threshold_antitone} proves $n/k\le n/m$ for every $n$ and every
$1\le m\le k$ using monotonicity of multiplication in $\mathbb N$.

An approximation order is represented by \texttt{UpperCut}, the set of
its rational upper bounds. For an ordinary numerical value $v$, this is
$\{q\in\mathbb Q_{\ge0}:v\le q\}$; it is upward closed, also when
$v=+\infty$ (in which case it is empty).
The formal theorem permits every upward-closed cut, a larger class than
real-valued approximation orders, so its universal impossibility conclusion
applies in particular to the usual numerical values.
\path{UpperCut.respects_equivalent} proves independence of the chosen
fraction representation.

To make the interpretation explicit, \texttt{OrderedValues} records any
numerical comparison relation with transitivity and an order-preserving interpretation
of rational bounds. \texttt{cutOfValue} constructs the corresponding
upper cut; \path{cutOfValue_exact} verifies that membership is precisely
the comparison $v\le q$. The final theorem \path{ordered_value_version}
states the contradiction directly using those numerical comparisons.
The project does not construct real numbers or the analytic definition
of $w_n$; neither is required, because the theorem is universal in the
values and uses only the indicated order laws.

\paragraph{Infinitude and source quantifiers.}
\texttt{InfinitelyOften p} means $\forall B\in\mathbb N,\ \exists n>B,\ p(n)$.
For a subset of $\mathbb N$ this is exactly infinitude. Passing to a weaker
inequality keeps each witness $n$, as verified by
\path{InfinitelyOften.mono}. \texttt{NoU} quantifies over every positive
integer index, and \texttt{W} uses the same family of excluded sets for
every $m$. The theorem \path{W_antitone} proves the full nestedness statement;
\path{conjecture388_refuted} then negates the two necessary clauses of
the original conjecture. Both this theorem and \path{ordered_value_version}
have no axiom dependencies. There are no added axioms, proof placeholders,
or native decision procedures. The project treats warnings as errors.
\end{document}
